\documentclass[12pt]{article}
\usepackage[margin=1in]{geometry}
\usepackage{amsmath, amssymb}
\usepackage{enumitem, array, xcolor}

\usepackage{boxedminipage}
\usepackage{fancyhdr, actuarialangle}
\usepackage{ragged2e, comment}


\newenvironment{problem}[1][]{%
  \bigskip% put space before problem statement box %
  \noindent\begin{boxedminipage}{\columnwidth}%
  \def\@tempa{#1}%
  \ifx\@tempa\empty\else%
    \textbf{#1}\hspace{0.5em}%
  \fi%
}{%
  \end{boxedminipage}%
}


\begin{document}

\begin{center}
    \Large \textbf{MATH 2790 Final Exam Fall 2025} \\[1em]
\end{center}

\vspace{.25cm}

\noindent Name: \underline{\hspace{8cm}} \hspace{.25cm} Auburn ID: \underline{\hspace{4cm}}



\section*{Instructions}
\begin{enumerate}[label=\arabic*.]
\item Print your full name and your banner ID on your exam. Then finish reading these instructions, and sign the bottom of this page.
\item All electronic equipment (cell phones, Apple Watch, AirPod, etc) are \textbf{NOT ALLOWED} during the exam time.
\item You are allowed to use up to 2 calculators from the following list: BA-35, BA II Plus, BA II Plus Professional, TI-30Xa, TI-30X II (IIS solar or IIB battery), TI-30XS MultiView (or XB battery).
\item No talking is allowed during the exam.
\item Unless otherwise indicated, \textbf{SHOW ALL YOUR WORK}. If no work is shown, no partial credit can be awarded. Your work and answer need to be accurate and relevant to receive points.
\item You will be given exactly 120 minutes for this exam.
\item Any student not following the above instructions nor behaving according to the above instructions during the exam may have their exam confiscated and points deducted.
\end{enumerate}

I have read and fully understand all of the above instructions. \\[2em]
Signature: \hrulefill


\vspace{.75cm}
Do not write in the table below

% scoring table
\begin{center}
\begin{tabular}{|*{3}{>{\centering\arraybackslash}p{.1\linewidth}|}}
    \hline
    \textbf{Problem} & \textbf{Points} & \textbf{Score} \\
    \hline
    1 & 11 &  \\
    \hline
    2 & 11 &  \\
    \hline
    3 & 11 &  \\
    \hline
    4 & 11 &  \\
    \hline
    5 & 11 &  \\
    \hline
    6 & 11 &  \\
    \hline
    7 & 11 &  \\
        \hline
    8 & 11 &  \\
        \hline
    9 & 11 &  \\
    \hline
    10 & 11 &  \\
    \hline \hline
    \textbf{Total} & 110 &  \\
    \hline
\end{tabular}
\end{center}

\footnotesize{Your grade will be scored out of 100. If you earned $X$ out of 110 points, you will receive an $X\%$.}





%%%%%%%%%%%%%%%%%%%%%%%%%%%%%%%%%%%%%%%%%%%%% 
\newpage 

\begin{problem}[1. (11 points)]
A 15-year bond is purchased at a premium. The bond pays semiannual coupons.
The amount of amortization of premium in the 6th coupon is 80.00. The amount of 
amortization of premium in the 27th coupon is 148.82. Calculate the premium of this bond. Round to 2 decimal places.
\end{problem}


%%%%%%%%%%%%%%%%%%%%%%%%%%%%%%%%%%%%%%%%%%%%% 

\newpage




\begin{problem}[2. (11 points)]
Given that the force of interest is $\delta_t = 0.02 + 0.01t$,
    \begin{enumerate}[label=(\alph*)]
        \item (9 points) Find the accumulation function $a(t)$.
        \item (6 points) Find the value at time $t=2$ of \$300 to be paid at time $t=5$. Round to 2 decimal places.
    \end{enumerate}
\end{problem}


%%%%%%%%%%%%%%%%%%%%%%%%%%%%%%%%%%%%%%%%%%%%% 
\newpage


\begin{problem}[3. (11 points)]
A small business owner, Maria, borrows \$25{,}000 from a local credit union to upgrade her equipment. Maria repays the loan by making a payment of \$10{,}000 at the end of one year, \$9{,}000 at the end of two years, and \$8{,}500 at the end of four years. The money received by the credit union is immediately reinvested at an annual effective rate of $2\%$ from years 1 to 2, and then at $3.5\%$ thereafter.
    \begin{enumerate}[label=(\alph*)]
        \item (8 points) Write out the equation to solve for Maria’s annual effective interest rate on the loan. Use the financial calculator to find the value. Round to 4 decimal places as a percentage.
        \item (8 points) Calculate the credit union's annual yield rate on this transaction. Round to 4 decimal places as a percentage.
    \end{enumerate}
\end{problem}

%%%%%%%%%%%%%%%%%%%%%%%%%%%%%%%%%%%%%%%%%%%%% 
\newpage



\begin{problem}[4. (11 points)]
An athlete, Emma, signs a sponsorship deal. She receives \$4 at the end of the first month, \$7 at the end of the second month, and thereafter each monthly payment increases by \$3. The amount of the last payment is \$64. 

\vspace{.25cm}

If for the first year, the nominal interest rate is 6\% convertible monthly and then increases to 9\% convertible monthly.  Calculate the value of Emma's sponsorship deal at the time of her first payment.
\end{problem}



%%%%%%%%%%%%%%%%%%%%%%%%%%%%%%%%%%%%%%%%%%%%% 
\newpage


\begin{problem}[5. (11 points)]
Lucy is awarded a scholarship that will pay her \$800 at the end of the first year. Each subsequent yearly payment increases by 1\%. The effective annual interest rate is 5\%.

\begin{enumerate}[label=(\alph*)]
    \item Suppose the scholarship payments continue for 10 years. Calculate the value of these payments at the time of the last payment.
    \item Now suppose instead that the scholarship payments continue forever. Lucy plans to reevaluate her finances at the time of her third payment. Calculate the value, at that time, of all scholarship payments (both past and future).
\end{enumerate}
\end{problem}


%%%%%%%%%%%%%%%%%%%%%%%%%%%%%%%%%%%%%%%%%%%%% 
\newpage

\begin{problem}[6. (11 points)]
A loan of $\$30{,}000$ is taken out at an annual effective interest rate of $2\%$.
Jordan makes annual payments that are five times the amount of interest accrued
in the year. He continues this process until he can make one final payment
of at most $\$250$. Find the time and the amount of Jordan's final payment.
Round to 2 decimal places.
\end{problem}

%%%%%%%%%%%%%%%%%%%%%%%%%%%%%%%%%%%%%%%%%%%%% 
\newpage


\begin{problem}[7. (11 points)]
For the next 15 years, Rosie will
deposit an amount of $\$1000$ at the end of each month into an account that earns a
nominal interest rate of $6\%$, compounded monthly. Immediately after her
last monthly deposit, the accumulated balance is transferred into a scholarship fund
that instead earns an annual effective interest rate of $4\%$. How large of a yearly
scholarship can Rosie give out forever, assuming that all of the deposited amount
goes toward the scholarship and the first scholarship is given one year after the account is opened?
\end{problem}

%%%%%%%%%%%%%%%%%%%%%%%%%%%%%%%%%%%%%%%%%%%%% 
\newpage


\begin{problem}[8. (11 points)]
Bond A is an $8$-year \$1{,}400 par-value bond with a nominal coupon rate of $8\%$,
convertible semiannually. 

\vspace{.25cm}

Bond B is a \$1{,}000 bond with a nominal coupon rate of $12\%$, convertible
semiannually, and a redemption amount of \$1{,}137.80. 

\vspace{.25cm}

If both bonds yield a nominal
rate of $6\%$ convertible semiannually and have the same price, how long (in years) until Bond B matures?
\end{problem}


%%%%%%%%%%%%%%%%%%%%%%%%%%%%%%%%%%%%%%%%%%%%% 
\newpage


\begin{problem}[9. (11 points)]
Account A is funded by an annuity-due with 10 payments of \$120 and earns an annual
effective interest rate of $5\%$. Account B is funded with a single deposit of \$1{,}200 today and earns simple interest at a rate of $3\%$ per year. Which account has a larger accumulated value at the end of 10 years? Justify your
answer.
\end{problem}



%%%%%%%%%%%%%%%%%%%%%%%%%%%%%%%%%%%%%%%%%%%%% 
\newpage



\begin{problem}[10. (11 points)]
A car is purchased for \$32{,}000. A down payment of \$2{,}000 is made at the time of
purchase, and the remaining balance is financed with a loan to be repaid with level
monthly payments over 6 years. The loan is subject to a nominal \textit{discount} rate
of $4.8\%$, convertible monthly.

\begin{enumerate}
    \item[(a)] Calculate the amount of the monthly payment. \textit{Round up} to 2 decimals.
    \item[(b)] Calculate the total interest paid during the first 3 years of the loan assuming the payments found in (a) are made every month excpet for a reduced final payment.
\end{enumerate}
\end{problem}











\newpage
Extra page for scratch Work

\newpage
Extra page for scratch Work



\end{document}
